\section{Reproducibility Package}
\label{app:repro}

\subsection{Artifact Structure}
The reproducibility package is organized into four logical components:
\begin{itemize}
  \item \textbf{Raw execution logs}: per-task, per-attempt records produced by each experimental mode, including prompts, candidate outputs, conformance scores, pattern hits, and timing.
  \item \textbf{Processed metrics}: aggregated per-task and per-mode summary tables derived from raw logs via a deterministic aggregation pipeline.
  \item \textbf{Figures}: publication-quality plots generated from processed metrics; each figure is associated with a manifest entry linking it to the corresponding data slice and research question.
  \item \textbf{Manuscript sources}: the \LaTeX{} source tree, including this document, tables, and figure references.
\end{itemize}

\subsection{Pipeline Stages}
Reproduction proceeds through four sequential stages:
\begin{enumerate}
  \item \textbf{Benchmark construction}: generate the hard benchmark tasks (180 tasks, up to 9 ordered scenarios each) using a deterministic task generator with a fixed seed.
  \item \textbf{Experiment execution}: run all compared modes (B0, B1, B2, M, A1, A2, A3) and the prevention evaluation protocol over the benchmark, producing raw logs.
  \item \textbf{Data preparation}: aggregate raw logs into per-mode metric summaries; this stage is idempotent given fixed raw logs and seed.
  \item \textbf{Figure and table generation}: produce all paper figures and statistical summary tables from the processed metrics.
\end{enumerate}
Each stage is independently re-runnable; pending or incomplete runs are marked \texttt{TBD} in reported tables rather than imputed.
All seeds, mode configurations, and stopping policies used in the paper are recorded in the artifact manifest for exact replication.
